\documentclass[10pt]{article}
\usepackage[a4paper, margin=0.75in]{geometry}
\usepackage{titlesec}
\usepackage{booktabs}
\usepackage{tabularx}
\usepackage{xcolor}
\usepackage{enumitem}
\usepackage{tgheros}
\usepackage[T1]{fontenc}
\usepackage{array}
\usepackage{textcomp}
\usepackage[hidelinks]{hyperref}
\usepackage{tikz}
\usetikzlibrary{shapes.rounded, shapes.geometric, arrows.meta, positioning, calc, fit, backgrounds}
\renewcommand{\familydefault}{\sfdefault}

%% ===== Colour palette =====
\definecolor{primary}{RGB}{15,52,96}        %% Deep petroleum navy
\definecolor{accent}{RGB}{198,124,0}        %% Oil amber
\definecolor{theme1}{RGB}{20,100,60}        %% Operations green
\definecolor{theme2}{RGB}{120,60,140}       %% Data violet
\definecolor{theme3}{RGB}{170,50,60}        %% Security red
\definecolor{theme4}{RGB}{30,90,140}        %% Cloud blue
\definecolor{lightgrey}{gray}{0.88}
\definecolor{midgrey}{gray}{0.65}
\definecolor{darkgrey}{gray}{0.38}

%% ===== Section style =====
\titleformat{\section}{%
  \color{primary}\bfseries\large\raggedright
}{}{0em}{}[\color{lightgrey}{\titlerule[1.5pt]}\vspace{-4pt}]

\titleformat{\subsection}{%
  \color{accent}\bfseries\normalsize\raggedright
}{}{0em}{}

\titlespacing*{\section}{0pt}{18pt}{8pt}
\titlespacing*{\subsection}{0pt}{12pt}{4pt}

\setlength{\parindent}{0pt}
\setlength{\parskip}{6pt}
\renewcommand{\arraystretch}{1.25}
\hypersetup{colorlinks=true, urlcolor=primary}

%% ===== Helpers =====
\newcommand{\spacer}{\vspace{10pt}}
\newcommand{\bigspacer}{\vspace{16pt}}

%% =====================================================
\begin{document}

%% ----- Title block -----
\begin{center}
  \vspace*{0.1cm}
  {\color{primary}\LARGE\bfseries NEWCROSS AND PAN OCEAN GROUP}\\[6pt]
  {\color{primary}\large\bfseries Digital Transformation --- A Perspective Paper}\\[10pt]
  {\color{darkgrey}\normalsize Version 2.0\quad\textbar\quad May 2026\quad\textbar\quad
   For Internal Discussion}
\end{center}

\bigspacer

%% ----- Contents -----
{\small
\renewcommand{\contentsname}{\color{primary}\normalsize\bfseries Contents}
\setcounter{tocdepth}{1}
\tableofcontents
}

\bigspacer

\noindent\colorbox{primary!8}{%
\begin{minipage}{\dimexpr\textwidth-2\fboxsep}
  \vspace{6pt}
  \textbf{\color{primary}About This Paper.}\quad
  This is a discussion paper, not a delivery plan. It is written from an external
  industry perspective and does not assume full visibility of initiatives already
  in motion inside the group. RADA-AMS, ongoing infrastructure modernisation, and
  other internal programmes are acknowledged as live work.

  \vspace{4pt}
  The objective here is to offer a frame --- the eight capability themes that
  define digital transformation in oil and gas in 2026, what they mean for an
  indigenous Nigerian operator, and where outside perspective and execution
  capacity could add value to what the group is already doing.
  \vspace{6pt}
\end{minipage}}

\bigspacer

%% ----------------------------------------------------------
\section{Why This Conversation, Why Now}
%% ----------------------------------------------------------

The Newcross and Pan Ocean Group of Companies already holds the most valuable
asset in any digital business: fifty years of operational data, deep domain
expertise, and a credibility inside the Nigerian oil and gas industry that no
foreign vendor can replicate.

The question is not whether to transform. The RADA-AMS commitment of October
2025 has already answered that. The question is what shape the transformation
takes, what capabilities the group invests in building, and how that work is
sequenced against operational reality.

This paper does not propose a project plan. It proposes a way of thinking about
the landscape, grounded in how the global industry is approaching the same
problem, and adapted to the group's specific context.

\bigspacer

%% ----------------------------------------------------------
\section{The Group of Companies}
%% ----------------------------------------------------------

The Newcross and Pan Ocean Group of Companies comprises three distinct legal
entities operating with partly shared infrastructure, shared senior leadership,
and complementary asset portfolios across the Nigerian upstream oil and gas
sector.

\spacer
\small
\begin{tabularx}{\columnwidth}{@{}lXl@{}}
\toprule
\textbf{Entity} & \textbf{Primary Role} & \textbf{Incorporated} \\
\midrule
Pan Ocean Oil Corporation (Nigeria) Ltd
  & Exploration, production, gas processing, and pipeline infrastructure.
    First indigenous company in JV with NNPC; OML-98 producing since 1976.
  & 1973 \\
Newcross Petroleum Limited
  & Upstream E\&P operations; corporate strategy and finance.
  & --- \\
Newcross Exploration and Production Limited
  & Dedicated exploration and drilling operations, including rig assets
    acquired during the Jonathan-era licensing round.
  & --- \\
\bottomrule
\end{tabularx}
\normalsize

\spacer
{\small\color{darkgrey}
\textbf{Leadership.}\enspace
Steven Fadeyi serves as Group Managing Director across all three entities.
Dr.\ Jason Fadeyi has focused on business strategy and finance across the
group. Digital systems designed for this group must support multi-entity
operation from day one --- not as a retrofit. The group is the customer,
not any single entity.
}

\bigspacer

%% ----------------------------------------------------------
\section{The Operational Portfolio}
%% ----------------------------------------------------------

\small
\begin{tabularx}{\columnwidth}{@{}lXl@{}}
\toprule
\textbf{Asset} & \textbf{Description} & \textbf{Status} \\
\midrule
OML-98 (Ogharefe field)
  & Original producing block; onshore northern Niger Delta. Operations
    commenced August 1976. Gas utilisation initiative running since 1984.
  & Producing \\
OML-147 (Northern Niger Delta)
  & Acquired as OPL-275 in 2007; PSC with NNPC signed 2009; converted to
    OML in 2015. Early Production Facility at Owa Aladinma:
    20{,}000 bopd plus 70 MMscf/d, 3-phase modularised design.
  & Producing \\
Ovade-Ogharefe Gas Plant (OOGP)
  & Cryogenic gas processing; 29 LPG storage tanks. Cornerstone of West
    Africa's largest carbon-emission reduction project. Phase 2 cryogenic
    expansion completed 2022.
  & Operating \\
Amukpe-Escravos Pipeline (AEPP)
  & 20-inch, 67 km, 160{,}000 bopd capacity. Built using Continuous
    Horizontal Directional Drilling (CHDD) anti-vandalism technology.
    Shared-use tariff arrangement with neighbouring producers.
  & Operating \\
Rig Assets (Newcross E\&P)
  & Drilling rig(s) acquired during the Jonathan-era licensing round, held
    under Newcross Exploration and Production Limited.
  & Active \\
\bottomrule
\end{tabularx}
\normalsize

\spacer
{\small\color{darkgrey}
\textbf{NNPC relationship.}\enspace
OML-98 was established under an original JV arrangement with NNPC --- that
structure has evolved significantly over the years. OML-147 operates under
a live PSC signed with NNPC in 2009; PSC reporting obligations on OML-147
remain active and are a direct driver of digital reporting requirements.
}

\bigspacer

%% ----------------------------------------------------------
\section{The Nigerian Oil and Gas Context}
%% ----------------------------------------------------------

Any serious transformation programme must be grounded in the regulatory and
market environment the group actually operates in. Several recent shifts
raise the urgency of the conversation.

\spacer
\small
\begin{tabularx}{\columnwidth}{@{}lX@{}}
\toprule
\textbf{Factor} & \textbf{Implication for the Group} \\
\midrule
Petroleum Industry Act 2021 (PIA)
  & Restructured fiscal regime; new royalty and tax framework; Host Community
    Development Trust (HCDT) obligation of 3\% of opex per entity. Tracking
    and reporting HCDT is a legal requirement, not optional. \\
NUPRC Regulatory Framework
  & Successor to DPR; expanded mandate, increased reporting cadence and data
    granularity. April 2026: directive mandating standardised templates for
    \textbf{measurement-based methane and GHG reporting} --- OOGP is directly
    in scope. \\
Nigerian Gas Flare Commercialisation Programme (NGFCP)
  & Pan Ocean's gas-flare leadership since 1984 is a credibility asset.
    Digital measurement and reporting of flare volumes is now an active
    compliance requirement. \\
Nigerian Content Act (NCA)
  & Local spend, employment, and training tracking across all three entities
    and both producing blocks. Audit risk is increasing. Currently manual
    throughout. \\
OML-147 PSC Reporting
  & Production, cost recovery, and fiscal data submitted to NNPC under PSC
    terms. Manually generated today from disparate sources. \\
Dangote Refinery Effect
  & Domestic crude market dynamics are shifting for the first time in
    decades. Oil marketing, lifting decisions, and buyer allocation now
    operate in a fundamentally different environment. \\
\bottomrule
\end{tabularx}
\normalsize

\bigspacer

%% ----------------------------------------------------------
\section{What Digital Transformation Means in Oil and Gas Today}
%% ----------------------------------------------------------

The phrase ``digital transformation'' is widely used and inconsistently defined.
When global majors and serious indigenous operators speak about it in 2026,
they are talking about building capability across eight themes. None of these
are projects. All of them are long-running programmes that mature over years.

\spacer

%% ===== Capability themes diagram (clean 4 x 2 grid + foundation bar) =====
\begin{center}
\begin{tikzpicture}[
  every node/.style={font=\small},
  theme/.style={draw, rounded corners=4pt,
                minimum width=3.7cm, minimum height=1.7cm,
                align=center, text=black, inner sep=4pt},
  found/.style={draw=primary, fill=primary!12, line width=0.8pt,
                rounded corners=5pt,
                minimum width=15.6cm, minimum height=1.0cm,
                align=center, font=\small\bfseries, text=primary,
                inner sep=6pt},
]

%% Row 1
\node[theme, fill=theme1!12, draw=theme1!70] (t1) at (0,0)
  {\textbf{1. Connected Operations}\\
   {\scriptsize Digital oilfield \textbullet\ SCADA \textbullet\ edge data}};
\node[theme, fill=theme1!12, draw=theme1!70, right=0.3cm of t1] (t2)
  {\textbf{2. Industrial AI}\\
   {\scriptsize Predictive maintenance \textbullet\ anomaly detection}};
\node[theme, fill=theme2!12, draw=theme2!70, right=0.3cm of t2] (t3)
  {\textbf{3. Data Fabric}\\
   {\scriptsize SAP \textbullet\ SCADA \textbullet\ geological \textbullet\ commercial}};
\node[theme, fill=theme2!12, draw=theme2!70, right=0.3cm of t3] (t4)
  {\textbf{4. ESG \& Decarbonisation}\\
   {\scriptsize Methane / GHG \textbullet\ emissions accounting}};

%% Row 2
\node[theme, fill=theme3!12, draw=theme3!70, below=0.35cm of t1] (t5)
  {\textbf{5. Workforce Digitisation}\\
   {\scriptsize Mobile ops \textbullet\ digital permits \textbullet\ knowledge}};
\node[theme, fill=theme3!12, draw=theme3!70, below=0.35cm of t2] (t6)
  {\textbf{6. OT \& Cyber Security}\\
   {\scriptsize Plant \textbullet\ pipeline \textbullet\ rig as attack surface}};
\node[theme, fill=theme4!12, draw=theme4!70, below=0.35cm of t3] (t7)
  {\textbf{7. Generative AI}\\
   {\scriptsize LLMs over well logs, contracts, reports}};
\node[theme, fill=theme4!12, draw=theme4!70, below=0.35cm of t4] (t8)
  {\textbf{8. Hybrid Cloud + Edge}\\
   {\scriptsize Offline-capable \textbullet\ edge inference}};

%% Foundation bar
\node[found, below=0.5cm of t6.south east, anchor=north]
  {Foundation\enspace\textbullet\enspace
   Connectivity\enspace\textbullet\enspace
   Identity \& Access\enspace\textbullet\enspace
   Data Governance};

\end{tikzpicture}
\end{center}

\spacer
{\small\color{darkgrey}
\textbf{Reading the diagram.}\enspace
The eight themes are not phases. They are concurrent capability areas that
mature at different rates. The shared foundation --- connectivity, identity,
and data governance --- is the boring infrastructure work without which none
of the eight themes can deliver value.
}

\bigspacer

%% ----------------------------------------------------------
\section{The Eight Themes --- A Group Perspective}
%% ----------------------------------------------------------

For each theme below: the industry definition, what it would look like if it
were mature at the group, and an honest note on where outside perspective is
likely to be useful versus where the group is probably already further along
than is visible from outside.

\spacer

\subsection*{1. Connected Operations}

\textbf{What it is.} Modernised SCADA, edge sensors, and reliable real-time
data flow from wells, plants, and pipelines into a usable data layer. The
foundation of every other theme.

\textbf{What mature looks like at the group.} OOGP, OML-98, OML-147, and the
AEPP corridor stream operational data into a common layer in near real time.
WAN, VSAT, and rig data links are monitored as data sources, not just IT
assets. Field decisions are made against current data, not yesterday's report.

\textbf{Note.} Significant work in this area is almost certainly already in
progress through SCADA programmes and infrastructure projects the external
view cannot see in detail. Outside perspective adds value mainly in the
integration layer and in treating connectivity health as a transformation
prerequisite, not a background IT task.

\spacer

\subsection*{2. Industrial AI and OT Analytics}

\textbf{What it is.} Machine learning applied to operational technology data
--- predictive maintenance, anomaly detection, production optimisation, and
asset integrity monitoring. Distinct from IT analytics: it operates on plant
and pipeline data, with the safety and reliability stakes that implies.

\textbf{What mature looks like at the group.} OOGP compressors, cryogenic
units, and LPG storage have failure-probability scoring before issues become
incidents. AEPP flow anomalies are detected automatically. Well performance
across OML-98 and OML-147 is continuously evaluated against historical
patterns.

\textbf{Note.} This is where AI delivers measurable economic value in oil
and gas. It depends on Theme 1 being in place first. The companion AI and
Data Roadmap document elaborates specific initiatives.

\spacer

\subsection*{3. Data Fabric and Single Source of Truth}

\textbf{What it is.} Breaking the silos between SAP, SCADA, geological systems,
commercial data, and HSE records. Treating data as an enterprise asset with
governance, lineage, ownership, and access control.

\textbf{What mature looks like at the group.} A single trusted view of
production, finance, compliance, and HSE across all three entities. Data
governance roles defined. Access controlled by entity and function. NDPR
compliance built in from the structure, not bolted on later.

\textbf{Note.} The hardest theme to demonstrate quickly and the most valuable
in the long run. Multi-entity data governance is genuinely difficult and is
where many transformation programmes lose momentum.

\spacer

\subsection*{4. ESG and Decarbonisation Digitisation}

\textbf{What it is.} Sensor-based measurement and digital reporting of methane,
GHG emissions, and other environmental metrics. Increasingly mandatory rather
than voluntary.

\textbf{What mature looks like at the group.} OOGP flare, methane, and process
emissions are measured continuously and feed directly into NUPRC-format
reports. Pan Ocean's gas-flare leadership since 1984 is reinforced by a
measurement-grade reporting capability.

\textbf{Note.} The April 2026 NUPRC directive on measurement-based methane
and GHG reporting makes this one urgent. It is not optional and the manual
fallback path is operationally untenable at OOGP scale.

\spacer

\subsection*{5. Workforce Digitisation}

\textbf{What it is.} Mobile-first field operations, digital permits-to-work,
incident and near-miss reporting, knowledge capture from experienced
engineers, immersive training where it adds value.

\textbf{What mature looks like at the group.} Field personnel at OOGP,
Ogharefe, Owa Aladinma, and rig sites work from mobile-accessible, offline-
capable systems. Experienced engineers' decision-making is captured before
they retire. New engineers learn from a digital institutional record, not
just oral tradition.

\textbf{Note.} HSE digital PTW is the most obvious starting point given the
high-consequence environment at OOGP. Knowledge capture from senior
engineering staff is a quietly urgent issue across the indigenous Nigerian
operator landscape.

\spacer

\subsection*{6. OT and Cyber-Physical Security}

\textbf{What it is.} The recognition that pipelines, plants, and rig control
systems are attack surfaces. OT security is now treated separately from IT
security and has become a board-level concern across the global industry in
the last three years.

\textbf{What mature looks like at the group.} OT environments are segmented
from IT. Plant, pipeline, and rig control systems have monitored access,
incident detection, and an explicit security posture. Vendor remote access is
controlled and logged.

\textbf{Note.} As more operational data flows out of plants and pipelines,
the attack surface grows. This theme is rarely the most visible internally
but is one of the highest-stakes areas to get right.

\spacer

\subsection*{7. Generative AI and Institutional Memory}

\textbf{What it is.} Large language models applied to decades of well logs,
drilling reports, geological surveys, contracts, and engineering documents.
Unlocking the institutional memory of long-producing assets.

\textbf{What mature looks like at the group.} An engineer can ask a natural-
language question against fifty years of OML-98 records and get a grounded
answer with sources. Contract terms, vendor history, and drilling rationale
are searchable rather than buried. Sensitive data stays inside the group's
infrastructure.

\textbf{Note.} OML-98's fifty-year archive is a strategic asset, not just a
records-management problem. Locally hosted models are essential here ---
well data and commercial history must not leave the group's infrastructure.

\spacer

\subsection*{8. Hybrid Cloud and Edge Architecture}

\textbf{What it is.} The recognition that pure cloud is the wrong architecture
for upstream oil and gas. Reliable operations require edge inference, offline-
capable systems, and cloud aggregation --- not cloud-only.

\textbf{What mature looks like at the group.} Field systems operate when WAN
links degrade. Critical inference runs locally at the asset. The cloud is the
aggregation, analytics, and consolidation layer --- not a single point of
failure for field operations.

\textbf{Note.} Foreign vendors selling cloud-only solutions for Nigerian
upstream operations are selling the wrong architecture. This is a place where
operational realism is itself a competitive advantage.

\bigspacer

%% ----------------------------------------------------------
\section{The Operational Complexity That Demands Change}
%% ----------------------------------------------------------

The eight themes above are abstract until grounded in what the group actually
runs. The complexity is real and is the primary source of transformation value.

\spacer
\small
\begin{tabularx}{\columnwidth}{@{}lX@{}}
\toprule
\textbf{Business Function} & \textbf{Current Pain and Opportunity} \\
\midrule
Exploration \& Production
  & Multi-block production data (OML-98 + OML-147) managed across separate
    systems; no unified real-time view; well performance managed manually. \\
Gas Plant Operations (OOGP)
  & Cryogenic equipment and LPG storage monitoring is reactive rather than
    predictive. Flare and GHG reporting is now a live NUPRC obligation
    (April 2026 directive). 29 LPG storage tanks make OOGP the group's
    highest-consequence asset for unplanned downtime. \\
Pipeline Management (AEPP)
  & 67 km corridor through the Niger Delta. Third-party injector metering
    and tariff reconciliation is paper-intensive. Crude theft and
    encroachment detection is reactive. Metering accuracy is a direct
    revenue issue. \\
Oil Marketing
  & Crude lifting schedules, buyer allocation, and PSC equity crude
    reconciliation managed manually. The Dangote refinery's entry into the
    domestic market is reshaping a system that was already stretched. \\
Finance \& Accounting
  & Multi-currency exposure (Naira + USD); intercompany transactions across
    three entities; NNPC PSC cost recovery accounting; annual budget cycle ---
    all heavily manual with significant reconciliation overhead. \\
Regulatory Reporting
  & NUPRC oil and gas returns; PSC reporting to NNPC; HCDT contribution
    tracking; Nigerian Content reports; GHG/methane reporting (fresh April
    2026 directive) --- all generated manually from disparate sources. \\
Health, Safety \& Environment
  & Permit-to-work systems paper-based. Incident reporting and near-miss
    tracking manual. OOGP is a high-consequence HSE environment by any
    definition (cryogenic processing; LPG storage). \\
HR \& Nigerian Content
  & Staff rotation scheduling, training compliance, and local spend
    classification for NCA reporting --- manual across three entities and
    multiple operational sites. \\
Procurement \& Supply Chain
  & Remote field operations with significant lead times; FX exposure on
    imported parts; emergency procurement at premium cost; vendor
    performance not systematically tracked. \\
IT and Communications Infrastructure
  & LAN, WAN, voice, and remote rig data links connect Lagos HQ, Ogharefe,
    Owa Aladinma, and rig sites. These links are not a background IT matter
    --- they are the physical foundation for every data initiative the group
    will pursue. A degraded WAN link to Ogharefe silences SCADA feeds and
    creates operational blind spots. \\
\bottomrule
\end{tabularx}
\normalsize

\bigspacer

%% ----------------------------------------------------------
\section{Capability Maturity --- A Self-Assessment Frame}
%% ----------------------------------------------------------

The matrix below is a frame for the group's own self-assessment, not an
external scoring. Mature transformation programmes review this kind of frame
quarterly. The honest answer for most operators sits in the middle columns
across most rows.

\spacer

{\small\color{darkgrey}
\textbf{How to use this matrix.}\enspace
Print the page. In a working session, mark one column per row based on the
group's own honest reading. The pattern of marks --- not any single row ---
is what reveals where investment is most needed and where the group is
already ahead of where outside observers would assume.
}

\spacer
\small
\begin{tabularx}{\columnwidth}{@{}lcccc@{}}
\toprule
\textbf{Capability Theme} & \textbf{Foundational} & \textbf{Developing} & \textbf{Operational} & \textbf{Differentiating} \\
\midrule
1. Connected Operations              & $\bigcirc$ & $\bigcirc$ & $\bigcirc$ & $\bigcirc$ \\
2. Industrial AI / OT Analytics      & $\bigcirc$ & $\bigcirc$ & $\bigcirc$ & $\bigcirc$ \\
3. Data Fabric                       & $\bigcirc$ & $\bigcirc$ & $\bigcirc$ & $\bigcirc$ \\
4. ESG and Decarbonisation           & $\bigcirc$ & $\bigcirc$ & $\bigcirc$ & $\bigcirc$ \\
5. Workforce Digitisation            & $\bigcirc$ & $\bigcirc$ & $\bigcirc$ & $\bigcirc$ \\
6. OT and Cyber-Physical Security    & $\bigcirc$ & $\bigcirc$ & $\bigcirc$ & $\bigcirc$ \\
7. Generative AI / Institutional Memory & $\bigcirc$ & $\bigcirc$ & $\bigcirc$ & $\bigcirc$ \\
8. Hybrid Cloud and Edge             & $\bigcirc$ & $\bigcirc$ & $\bigcirc$ & $\bigcirc$ \\
\bottomrule
\end{tabularx}
\normalsize

\spacer
{\small
\textbf{Levels.}\quad
\textbf{Foundational:} basic infrastructure exists. \quad
\textbf{Developing:} pilots running, value not yet measurable. \quad
\textbf{Operational:} embedded in business as usual, measurable value. \quad
\textbf{Differentiating:} a competitive advantage in the Nigerian market.
}

\bigspacer

%% ----------------------------------------------------------
\section{The Compliance Imperative}
%% ----------------------------------------------------------

Several capability investments are not discretionary choices --- they are
responses to legal obligations already in force or recently activated. This
reframes the conversation from ``efficiency project'' to ``compliance
infrastructure.''

\spacer

\noindent\colorbox{primary!8}{%
\begin{minipage}{\dimexpr\textwidth-2\fboxsep}
  \vspace{6pt}
  \textbf{\color{primary}April 2026 --- NUPRC Methane and GHG Directive.}\quad
  NUPRC issued a directive mandating standardised templates and the transition
  to \textbf{measurement-based methane and GHG reporting} across all upstream
  operations. OOGP is directly in scope. Manual compliance with this directive
  is operationally untenable at scale. An automated sensor-to-report pipeline
  is the only sustainable response and sits inside Theme 4.
  \vspace{6pt}
\end{minipage}}

\spacer
\small
\begin{tabularx}{\columnwidth}{@{}lXl@{}}
\toprule
\textbf{Obligation} & \textbf{Scope} & \textbf{Status} \\
\midrule
HCDT Contribution Tracking
  & 3\% of opex per PIA 2021; applies across all three group entities
  & In force \\
NUPRC Methane / GHG Reporting
  & Measurement-based; standardised templates; OOGP directly in scope
  & Fresh, April 2026 \\
Nigerian Content Act Reporting
  & Local spend, employment, and training across all entities and blocks
  & Ongoing \\
OML-147 PSC Reporting to NNPC
  & Production, cost recovery, and fiscal data under live PSC terms
  & Ongoing \\
Gas Flare Reporting (NGFCP)
  & Flare volume measurement and programme participation
  & Ongoing \\
NDPR Compliance
  & Nigeria Data Protection Regulation applies to employee data and
    some operational datasets
  & Ongoing \\
\bottomrule
\end{tabularx}
\normalsize

\bigspacer

%% ----------------------------------------------------------
\section{Where External Perspective Adds Value}
%% ----------------------------------------------------------

This paper does not propose to redirect work that is already in motion.
Internal teams have visibility, context, and relationships that no external
perspective can substitute for. The honest contribution of an outside view is
narrower:

\spacer
\small
\begin{tabularx}{\columnwidth}{@{}lX@{}}
\toprule
\textbf{Contribution} & \textbf{Why It Matters} \\
\midrule
Industry pattern recognition
  & How global majors and other indigenous operators are sequencing the
    eight themes; what has worked and what has consistently failed. \\
Capability framing
  & Treating transformation as concurrent capability areas rather than a
    linear programme. This frame is what most internal plans lack, not for
    want of effort but because internal momentum tends to push toward
    project plans. \\
Architecture discipline
  & Multi-entity, offline-first, edge-aware architecture from day one.
    Single-entity designs and cloud-only assumptions are expensive
    mistakes that show up only in year two. \\
Commercialisation thinking
  & The same capabilities that solve the group's problems are sellable
    products in the Nigerian market. Designing for that possibility from
    the start costs little; retrofitting later is rarely possible. \\
Independent technical review
  & Outside engineering perspective on architectural choices, vendor
    evaluation, and risk assessment --- complementary to internal teams,
    not competitive with them. \\
\bottomrule
\end{tabularx}
\normalsize

\bigspacer

%% ----------------------------------------------------------
\section{Building for Industry as Well as Internal Use}
%% ----------------------------------------------------------

The capabilities the group builds for itself are, with few exceptions, the
same capabilities every Nigerian operator needs. Designing for that reality
from the start is a discipline question, not a future decision. It costs
little upfront. Retrofitting it later is rarely possible.

\spacer

\noindent\colorbox{primary!8}{%
\begin{minipage}{\dimexpr\textwidth-2\fboxsep}
  \vspace{6pt}
  \textbf{\color{primary}The Discipline.}\quad
  Build once for the group. Design for the industry from day one. Configurable
  multi-entity support, clean APIs, and tenant-ready data models are the
  difference between a bespoke IT project and a product. The group's fifty
  years of operating history is the most powerful reference deployment any
  Nigerian software business could hold.
  \vspace{6pt}
\end{minipage}}

\spacer

Candidate capabilities with clear external market potential, drawn from the
companion AI and Data Roadmap:

\spacer
\small
\begin{tabularx}{\columnwidth}{@{}lX@{}}
\toprule
\textbf{Capability} & \textbf{Why It Travels Beyond the Group} \\
\midrule
Regulatory reporting automation
  & NUPRC, NCA, HCDT, and PSC reporting are mandatory for every upstream
    operator. Foreign vendors consistently fail on Nigerian regulatory
    specifics. \\
Predictive maintenance for Nigerian gas plants
  & Every gas plant operator in Nigeria has the same problem. OOGP is a
    credible reference site. \\
Pipeline integrity and tariff reconciliation
  & Niger Delta operators with long-haul pipelines face identical
    challenges. AEPP is a 67 km reference site. \\
Multi-entity operations dashboard
  & The Nigerian indigenous sector has consolidated significantly; multi-
    entity operations are now the norm, not the exception. \\
HSE digital systems for Nigerian site conditions
  & Foreign HSE software does not understand Nigerian operational reality
    or NUPRC reporting requirements. \\
\bottomrule
\end{tabularx}
\normalsize

\bigspacer

%% ----------------------------------------------------------
\section{Capability and Team Considerations}
%% ----------------------------------------------------------

Without commenting on roles already in place, the patterns below recur in
oil and gas transformation programmes that succeed.

\spacer
\small
\begin{tabularx}{\columnwidth}{@{}lX@{}}
\toprule
\textbf{Pattern} & \textbf{Why It Matters} \\
\midrule
Domain-fluent leadership
  & The lead must understand both the operational business and modern
    software practice. Hiring exclusively from one side or the other
    consistently underdelivers. \\
Internal + external mix
  & Internal hires bring institutional trust and domain knowledge.
    External engineers bring modern practice. Both are required;
    neither alone is sufficient. \\
Diaspora engineering talent
  & Engineers with Nigerian context and global technical exposure are
    a structural advantage that few competitors can match. \\
DevOps and platform discipline early
  & The single largest predictor of whether internal solutions can become
    products. Underinvestment here is the most common transformation
    failure pattern. \\
Business analyst rooted in operations
  & The hardest role to fill externally and the most important to get
    right. Without operational fluency at the requirements layer,
    technically capable teams build the wrong thing. \\
\bottomrule
\end{tabularx}
\normalsize

\bigspacer

%% ----------------------------------------------------------
\section{Common Failure Modes}
%% ----------------------------------------------------------

\small
\begin{tabularx}{\columnwidth}{@{}lXX@{}}
\toprule
\textbf{Risk} & \textbf{What Happens} & \textbf{Counter} \\
\midrule
Building for presentations
  & Impressive demos that do not survive contact with real operations
  & Ship working software early; iterate from actual field usage \\
Ignoring infrastructure constraints
  & Solutions requiring constant internet at Ogharefe, Owa Aladinma, or
    along the AEPP corridor fail immediately
  & Offline-first architecture from day one \\
Over-relying on external advisors without operational grounding
  & Limited understanding of NUPRC, the Niger Delta, Naira FX, and Nigerian
    site conditions; recommendations look reasonable on paper but fail on
    contact with operational reality
  & Build internal capability first; engage external advisors for specific
    technical skills and pattern recognition, not for strategy ownership \\
Single-entity design
  & Systems that cannot accommodate all three entities; commercialisation
    becomes impossible
  & Multi-entity architecture from day one; the group is the customer \\
Bespoke instead of product
  & Every solution is unique; nothing reusable across entities or sellable
    externally
  & API-first; configuration over customisation from the start \\
Treating transformation as a project
  & Time-boxed plans collapse on contact with operational reality;
    momentum is lost
  & Capability-themed programme; year-on-year maturity, not phase exits \\
Ignoring data quality
  & AI trained on bad data produces bad and potentially dangerous
    decisions
  & Data governance and audit before any ML project begins \\
\bottomrule
\end{tabularx}
\normalsize

\bigspacer

%% ----------------------------------------------------------
\section{Guiding Principles}
%% ----------------------------------------------------------

\small
\begin{tabularx}{\columnwidth}{@{}lX@{}}
\toprule
\textbf{Principle} & \textbf{What It Means in Practice} \\
\midrule
Nigerian problems first
  & Build for the NUPRC, the Niger Delta, the Naira, and the generator.
    Not for assumptions about infrastructure that do not hold here. \\
Offline-first
  & Core functions must operate without reliable internet at every
    operational site. \\
Explainable AI
  & Operations managers must understand why the system is recommending
    something. Black-box AI will not be adopted by field teams. \\
Data before models
  & Data audit and governance precedes any machine learning initiative.
    Dirty data produces dangerous recommendations. \\
Locally hosted where sensitive
  & Well data, production history, and commercial data remain on-premise
    or within the group's tenant. Not on third-party cloud APIs. \\
Multi-entity from day one
  & Single-entity designs cannot be retrofitted later. The group is the
    customer; design accordingly. \\
Capability over project
  & Long-running themes that mature over years, not phases that end and
    are declared complete. \\
RADA-AMS is the mandate
  & The group publicly committed to energy and innovation leadership through
    RADA-AMS in October 2025. This work operationalises that commitment ---
    it does not replace the existing strategic direction. \\
\bottomrule
\end{tabularx}
\normalsize

\bigspacer

%% ----- About the Author -----
\noindent\colorbox{primary!8}{%
\begin{minipage}{\dimexpr\textwidth-2\fboxsep}
  \vspace{6pt}
  \textbf{\color{primary}About the Author.}\quad
  Olajide Olaniyan is a Chartered Engineer (CEng) and a software engineer
  at IBM, working on Cloud Pak for AIOps --- IBM's enterprise platform for
  AI-driven IT operations and observability. He previously supported the
  Newcross and Pan Ocean Group's IT and communications infrastructure
  between 2010 and 2021, with operational responsibility for the group's
  LAN, WAN, and remote site connectivity across Lagos, Ogharefe, and
  associated operational sites.

  \vspace{4pt}
  This paper draws on that prior operational grounding combined with
  current exposure to enterprise-scale AI and data platforms. It does not
  assume continued visibility into the group's internal systems and is
  written deliberately at the level of industry perspective and capability
  framing rather than implementation detail.

  \vspace{4pt}
  This paper is shared in a personal capacity. Any subsequent engagement
  would be delivered through Atensai --- an Ireland-based technology
  advisory and engineering practice currently in formation, with a
  Nigerian arm under discussion.
  \vspace{6pt}
\end{minipage}}

\bigspacer

%% ----- Footer -----
{\small\color{darkgrey}
\textbf{Newcross and Pan Ocean Group of Companies}\enspace|\enspace
Strategic Vision --- Version 2.0\enspace|\enspace May 2026\enspace|\enspace
Perspective Paper\\
Companion document: \textit{AI and Data Roadmap --- Newcross and Pan Ocean
Group}.\\
This paper offers an external industry perspective and does not assume full
visibility of internal initiatives already in motion.
}

\end{document}
